% GENERATED by scripts/build-lecture-slides.py; do not edit. Portable profile: XeTeX, theme and fonts from the core.
\makeatletter\IfFileExists{core/beamer/beamerthememx.sty}{\def\mxcorepath{core/}\def\input@path{{core/beamer/}{core/}}}{\def\mxcorepath{../core/}\def\input@path{{../core/beamer/}{../core/}}}\makeatother
\documentclass[aspectratio=169,10pt,t]{beamer}
\usetheme{mx}
\usepackage{graphicx,booktabs}
\usepackage[french]{babel}
\setlength{\parskip}{.45em}
\hypersetup{pdfauthor={},pdfsubject={Lecture slides},pdfkeywords={},pdftitle={Suites numériques}}
\providecommand{\N}{\mathbb{N}}
\providecommand{\Z}{\mathbb{Z}}
\providecommand{\D}{\mathbb{D}}
\providecommand{\Q}{\mathbb{Q}}
\providecommand{\R}{\mathbb{R}}
\providecommand{\C}{\mathbb{C}}
\providecommand{\vect}[1]{\overrightarrow{#1}}
\providecommand{\vi}{\vec{\imath}}
\providecommand{\vj}{\vec{\jmath}}
\providecommand{\vk}{\vec{k}}
\providecommand{\norm}[1]{\left\lVert #1\right\rVert}
\providecommand{\vc}[2]{\begin{pmatrix}#1\\#2\end{pmatrix}}
\providecommand{\deter}[2]{\det\left(#1,#2\right)}
\providecommand{\abs}[1]{\left\lvert #1\right\rvert}
\providecommand{\dd}{\,\mathrm{d}}
\providecommand{\e}{\mathrm{e}}
\providecommand{\Card}{\operatorname{Card}}
\providecommand{\E}{\operatorname{E}}
\providecommand{\V}{\operatorname{V}}
\providecommand{\degre}{^{\circ}}
\providecommand{\un}[1]{\,\mathrm{#1}}
\providecommand{\sci}[2]{#1\times10^{#2}}
\providecommand{\pgcd}{\operatorname{PGCD}}
\providecommand{\ppcm}{\operatorname{PPCM}}
\title{Suites numériques}
\subtitle{Cours et exemples}
\date{}
\author{}
\begin{document}
\mxtitleframe
\begin{frame}{Objectifs}
Définir et représenter une suite, déterminer ses variations et sa limite, utiliser les suites arithmétiques et géométriques et calculer leurs sommes.
\end{frame}
\begin{frame}{Une fonction des entiers naturels}
Une suite réelle \((u_n)\) est une fonction définie sur \(\mathbb{N}\) (ou sur les entiers à partir d’un rang précisé). Une formule explicite donne \(u_n\) directement; une relation de récurrence donne le terme suivant à partir des précédents et nécessite des termes initiaux. On représente les points \((n;u_n)\) : seuls les indices entiers appartiennent au domaine.
\end{frame}
\begin{frame}{Une fonction des entiers naturels}
\(u_n=n^2+1\) donne \(1,2,5,10\) aux rangs \(0,1,2,3\). Avec \(v_0=1\) et \(v_{n+1}=2v_n+1\), on obtient \(1,3,7,15\). Ce sont deux modes de génération différents.
\end{frame}
\begin{frame}{Une fonction des entiers naturels}
Une suite est croissante si \(u_{n+1}\ge u_n\) pour chaque rang; décroissante avec l’inégalité opposée. Elle converge vers \(\ell\) si ses termes se rapprochent de \(\ell\) autant que voulu à partir d’un rang suffisamment grand. Sinon elle est divergente; une divergence peut être infinie ou oscillante.
\end{frame}
\begin{frame}{Une fonction des entiers naturels}
Calculer \(u_{n+1}-u_n\). Pour des termes strictement positifs, comparer aussi \(u_{n+1}/u_n\) à \(1\). Un quotient n’est pas utilisable si un terme est nul; la positivité est nécessaire pour conserver le sens de l’inégalité.
\end{frame}
\begin{frame}{Suites arithmétiques}
\(u_{n+1}=u_n+r\) pour un réel constant \(r\) définit une suite arithmétique de raison \(r\).
\end{frame}
\begin{frame}{Suites arithmétiques}
Pour \(n\ge p\), \(u_n=u_p+(n-p)r\). La somme de \(k\ge1\) termes à partir de \(u_p\) vaut
\[
 u_p+\cdots+u_{p+k-1}=\dfrac{k(u_p+u_{p+k-1})}{2}.
\]
\end{frame}
\begin{frame}{Suites arithmétiques}
Ajouter \(r\) exactement \(n-p\) fois donne le terme général, formalisable par récurrence. Pour la somme, écrire les termes dans les deux ordres et additionner : chaque paire vaut premier plus dernier; deux fois la somme vaut \(k\) fois cette quantité.
\end{frame}
\begin{frame}{Suites géométriques}
\(u_{n+1}=qu_n\), pour un réel constant \(q\), définit une suite géométrique. Pour \(q=0\), tous les termes après le terme initial sont nuls.
\end{frame}
\begin{frame}{Suites géométriques}
Pour \(q\ne0\), \(u_n=u_pq^{n-p}\). Pour \(q=0\), cette formule vaut pour \(n>p\) et \(u_p\) reste le terme initial, sans utiliser \(0^0\). Si \(q\ne1\),
\[
 u_p+\cdots+u_{p+k-1}=u_p\dfrac{1-q^k}{1-q}.
\]
Si \(q=1\), la somme vaut \(ku_p\). La formule de somme convient aussi à \(q=0\).
\end{frame}
\begin{frame}{Suites géométriques}
Le terme général résulte de multiplications successives. Pour \(q=0\), la somme de \(k\ge1\) termes vaut \(u_p\), ce qui coïncide avec la formule annoncée. Pour \(q\ne0\), considérer la somme \(S=u_p+u_pq+\cdots+u_pq^{k-1}\), soustraire \(qS\) donne \((1-q)S=u_p(1-q^k)\). Séparer \(q=1\) avant la division.
\end{frame}
\begin{frame}{Suites géométriques}
Si \(|q|<1\), \(q^n\to0\). Si \(q>1\), \(q^n\to+\infty\). Si \(q=-1\), les termes alternent; si \(q<-1\), leur module croît sans limite et leurs signes alternent. Ainsi, pour une suite géométrique non nulle, la convergence dépend de ces cas; pour \(q=1\), la suite est constante. Si \(u_0=0\), elle est nulle quel que soit \(q\).
\end{frame}
\begin{frame}{Un modèle de récurrence}
Si \(u_{n+1}=0{,}8u_n+10\), la valeur fixe vérifie \(L=0{,}8L+10\), donc \(L=50\). Poser \(v_n=u_n-50\) donne \(v_{n+1}=0{,}8v_n\), d’où \(u_n=50+(u_0-50)(0{,}8)^n\) et \(u_n\to50\).
\end{frame}
\begin{frame}{Un modèle de récurrence}
Pour \(u_{n+1}=au_n+b\), si \(a\ne1\), la valeur fixe est \(b/(1-a)\). Si \(a=1\), la suite est arithmétique; il ne faut pas diviser par \(1-a\). Ce changement de suite est une méthode d’exercice, pas une nouvelle formule à apprendre sans preuve.
\end{frame}
\begin{frame}{Un modèle de récurrence}
De \(u_p\) à \(u_q\) inclus, il y a \(q-p+1\) termes, et non \(q-p\). Identifier la nature de la suite et ses paramètres avant de choisir la formule. Si les termes se composent d’une constante et d’une géométrique, sommer séparément ces deux parties.
\end{frame}
\begin{frame}{Un modèle de récurrence}
Pour \(u_n=10+6(1/2)^n\), la somme \(u_0+\cdots+u_3\) comporte quatre termes. La partie constante vaut \(40\); la partie géométrique vaut \(6(1-(1/2)^4)/(1-1/2)=45/4\). Au total \(S=205/4=51{,}25\). Contrôle direct : \(16+13+11{,}5+10{,}75\) donne le même résultat.
\end{frame}
\begin{frame}{Un modèle de récurrence}
Pour \(u_{n+1}=au_n+b\) avec \(a\ne1\), résoudre \(L=aL+b\), puis calculer réellement \(v_{n+1}=u_{n+1}-L\) pour obtenir \(av_n\). Cette identité, avec le terme initial, donne la formule. Étudier ensuite \(u_{n+1}-u_n\) en conservant le signe de \(u_0-L\) : une raison entre \(0\) et \(1\) ne rend pas automatiquement la suite décroissante.
\end{frame}
\begin{frame}{Un modèle de récurrence}
Avec \(u_{n+1}=u_n/2+5\), \(L=10\) et \(u_n=10+(u_0-10)(1/2)^n\). Si \(u_0>10\), les termes sont au-dessus de \(10\) et décroissent; si \(u_0<10\), ils sont en dessous et croissent. Si \(u_0=10\), la suite est constante. Dans les trois cas, la limite est \(10\).
\end{frame}
\begin{frame}{Un modèle de récurrence}
Le terme \(u_0\) est le premier terme d’une suite commençant à zéro. Relier les points pour guider l’œil ne transforme pas le domaine discret en intervalle réel. Dans une question de premier rang, vérifier le rang précédent pour établir la minimalité.
\end{frame}
\end{document}
